\documentclass{beamer}

\usetheme{Berlin}
\usecolortheme{Orchid}
\useoutertheme{miniframes}
\setbeamertemplate{navigation symbols}{}

\usepackage{amsmath}
\usepackage{booktabs}
\usepackage{graphicx}
\usepackage{xcolor}
\usepackage{listings}
\usepackage{hyperref}

\lstset{
  basicstyle=\ttfamily\small,
  keywordstyle=\color{blue},
  commentstyle=\color{gray},
  stringstyle=\color{teal},
  showstringspaces=false
}

\title[Python Programming]{Python Programming}
\subtitle{Unit 02 -- Lecture 03: Sets and Dictionaries}
\author{Tofik Ali}
\institute{School of Computer Science, UPES Dehradun}
\date{\today}

\begin{document}

\begin{frame}[fragile]
  \titlepage
  \vspace{0.5em}
  \begin{center}
  \footnotesize Repository: \texttt{https://github.com/tali7c/Python-Programming}
  \end{center}
\end{frame}

\begin{frame}[fragile]{Quick Links}
  \centering
  \hyperlink{sec:core}{\beamerbutton{Core Concepts}}\hspace{1em}
  \hyperlink{sec:demo}{\beamerbutton{Demo}}\hspace{1em}
  \hyperlink{sec:interactive}{\beamerbutton{Interactive}}\hspace{1em}
  \hyperlink{sec:summary}{\beamerbutton{Summary}}
\end{frame}

\begin{frame}[fragile]{Agenda}
  \tableofcontents
\end{frame}

\section{Overview}

\section{Core Concepts}
\label{sec:core}

\begin{frame}[fragile]{Learning Outcomes}
  \begin{itemize}[<+->]
    \item Explain why sets guarantee uniqueness and why they do not support indexing
    \item Apply set methods (\texttt{add/update/remove/discard}) and set operations (union, intersection, etc.)
    \item Build dictionaries for real-world key--value modeling and safe lookup
    \item Work with nested dictionaries/lists and sort dictionary items by key/value
    \item Use typecasting (list/tuple/set/dict) for data cleaning and restructuring
  \end{itemize}
\end{frame}

% ----------------------
% Sets
% ----------------------

\begin{frame}[fragile]{Sets: Definition and Properties}
  \begin{itemize}[<+->]
    \item Unordered collection of \textbf{unique} elements
    \item Fast membership testing: \texttt{x in s}
    \item Mutable (you can add/remove elements)
    \item Elements must be \textbf{hashable} (e.g., int/str/tuple, not list)
  \end{itemize}
\end{frame}

\begin{frame}[fragile]{Creating Sets (Including Empty Sets)}
  \begin{itemize}[<+->]
    \item Use braces for non-empty sets: \texttt{\{1, 2, 3\}}
    \item Use \texttt{set(...)} to convert from any iterable
    \item \texttt{\{\}} is an empty dictionary, so use \texttt{set()} for an empty set
  \end{itemize}
  \vspace{0.3em}
  \begin{lstlisting}[language=Python]
A = {1, 2, 2, 3}          # duplicates removed
B = set([3, 4, 4, 5])     # from list
empty = set()             # NOT {}
  \end{lstlisting}
\end{frame}

\begin{frame}[fragile]{Set Methods (Add/Remove/Clear/Copy)}
  \small
  \begin{tabular}{ll}
    \toprule
    Method & What it does \\
    \midrule
    \texttt{add(x)} & add one element \\
    \texttt{update(iter)} & add many elements \\
    \texttt{remove(x)} & remove; error if missing \\
    \texttt{discard(x)} & remove; no error if missing \\
    \texttt{pop()} & remove and return an arbitrary element \\
    \texttt{clear()} & remove all elements \\
    \texttt{copy()} & shallow copy \\
    \bottomrule
  \end{tabular}
  \normalsize
\end{frame}

\begin{frame}[fragile]{Set Operations: Operators vs Methods}
  \begin{itemize}[<+->]
    \item Union: \texttt{A | B} \hspace{1em} or \hspace{1em} \texttt{A.union(B)}
    \item Intersection: \texttt{A \& B} \hspace{0.75em} or \hspace{0.75em} \texttt{A.intersection(B)}
    \item Difference: \texttt{A - B} \hspace{1em} or \hspace{1em} \texttt{A.difference(B)}
    \item Symmetric difference: \texttt{A \string^ B} \hspace{0.35em} or \hspace{0.35em} \texttt{A.symmetric\_difference(B)}
  \end{itemize}
  \vspace{0.3em}
  \begin{lstlisting}[language=Python]
A = {1, 2, 3}
B = {2, 3, 4}
print(A & B)   # {2, 3}
  \end{lstlisting}
\end{frame}

\begin{frame}[fragile]{Set Comparisons (Subset/Superset/Disjoint)}
  \begin{itemize}[<+->]
    \item Subset: \texttt{A <= B} or \texttt{A.issubset(B)}
    \item Superset: \texttt{A >= B} or \texttt{A.issuperset(B)}
    \item Disjoint: \texttt{A.isdisjoint(B)} (no common elements)
  \end{itemize}
\end{frame}

% ----------------------
% Dictionaries
% ----------------------

\begin{frame}[fragile]{Dictionaries: Key--Value Model}
  \begin{itemize}[<+->]
    \item Map keys to values (fast lookup by key)
    \item Keys must be unique and hashable (immutable)
    \item Values can be any type (including lists/dicts)
    \item In modern Python, dictionaries preserve insertion order (but think ``mapping'')
  \end{itemize}
  \vspace{0.3em}
  \begin{lstlisting}[language=Python]
student = {"name": "Asha", "marks": 88}
print(student["name"])
  \end{lstlisting}
\end{frame}

\begin{frame}[fragile]{Creating Dictionaries (Common Ways)}
  \begin{itemize}[<+->]
    \item Curly braces: \texttt{\{"a": 1, "b": 2\}}
    \item \texttt{dict()} constructor: \texttt{dict(a=1, b=2)}
    \item From a list of pairs: \texttt{dict([("a", 1), ("b", 2)])}
  \end{itemize}
\end{frame}

\begin{frame}[fragile]{Essential Dictionary Methods}
  \begin{itemize}[<+->]
    \item Views: \texttt{d.keys()}, \texttt{d.values()}, \texttt{d.items()}
    \item Safe access: \texttt{d.get(key, default)}
    \item Update/merge: \texttt{d.update(other)}
    \item Removal: \texttt{d.pop(key)}, \texttt{d.popitem()}
    \item Cleanup: \texttt{d.clear()}
  \end{itemize}
  \vspace{0.3em}
  \begin{lstlisting}[language=Python]
freq = {}
word = "python"
freq[word] = freq.get(word, 0) + 1
  \end{lstlisting}
\end{frame}

\begin{frame}[fragile]{Nested Dictionaries and Lists}
  \begin{itemize}[<+->]
    \item Real data is often nested (JSON-like)
    \item Access step-by-step using keys and indices
  \end{itemize}
  \vspace{0.3em}
  \begin{lstlisting}[language=Python]
response = {"user": {"profile": {"email": "k@gmail.com"}}}
email = response["user"]["profile"]["email"]
  \end{lstlisting}
\end{frame}

\begin{frame}[fragile]{Sorting Dictionary Data}
  \begin{itemize}[<+->]
    \item Sort by key: \texttt{sorted(d.items())}
    \item Sort by value: \texttt{sorted(d.items(), key=lambda kv: kv[1])}
    \item Ranking (top-k): sort then slice
  \end{itemize}
  \vspace{0.3em}
  \begin{lstlisting}[language=Python]
counts = {"a": 3, "b": 1, "c": 2}
top = sorted(counts.items(), key=lambda kv: kv[1], reverse=True)[:2]
  \end{lstlisting}
\end{frame}

\begin{frame}[fragile]{Typecasting Between Collections (Data Cleaning)}
  \begin{itemize}[<+->]
    \item Remove duplicates: \texttt{list(set(data))} (order may change)
    \item Modify an immutable tuple: \texttt{tuple -> list -> tuple}
    \item Build dict from pairs: \texttt{dict([("a", 1), ("b", 2)])}
  \end{itemize}
\end{frame}

\begin{frame}[fragile]{Where Sets and Dictionaries Shine}
  \begin{itemize}[<+->]
    \item Sets: unique email addresses, fast membership checks, comparing groups
    \item Dictionaries: student records, inventory, shopping carts, frequency counts
    \item Both: building reliable, structured data models for programs
  \end{itemize}
\end{frame}

\section{Demo}
\label{sec:demo}

\begin{frame}[fragile]{Demo: Word Frequency and Unique Words}
  \begin{itemize}[<+->]
    \item Split a sentence into words
    \item Build a frequency dictionary (\texttt{get(..., 0)} pattern)
    \item Build a set of unique words
    \item Extension: print top 3 most frequent words
  \end{itemize}
  \vspace{0.4em}
  \textbf{Script:} \texttt{demo/word\_frequency.py}
\end{frame}

\section{Interactive}
\label{sec:interactive}

\begin{frame}[fragile]{Checkpoint 1}
  Which code correctly creates an empty set?
  \begin{enumerate}[<+->]
    \item \texttt{s = \{\}}
    \item \texttt{s = set()}
  \end{enumerate}
\end{frame}

\begin{frame}[fragile]{Checkpoint 2}
  What is the difference between \texttt{s.remove(x)} and \texttt{s.discard(x)}?
\end{frame}

\begin{frame}[fragile]{Checkpoint 3}
  Given \texttt{response = \{"user": \{"profile": \{"email": "k@gmail.com"\}\}\}},\\
  write the expression to access the email value.
\end{frame}

\begin{frame}[fragile]{Think-Pair-Share}
  Choose the right collection type (list/tuple/set/dict) for:
  \begin{itemize}[<+->]
    \item unique roll numbers in a class
    \item a fixed (x, y) coordinate
    \item a student's profile (name, roll, marks)
  \end{itemize}
\end{frame}

\section{Summary}
\label{sec:summary}

\begin{frame}[fragile]{Summary}
  \begin{itemize}[<+->]
    \item Sets store unique, unordered values and support set algebra
    \item Dictionaries map keys to values for fast lookup and counting
    \item Nested structures are common in real data (JSON-like)
    \item Sorting and typecasting help with analytics and data cleaning
  \end{itemize}
\end{frame}

\begin{frame}[fragile]{Exit Question}
  Write one line to sort dictionary items by value in descending order.
\end{frame}

\end{document}

